\section{The Algorithm Toolbox: Three Axes}
\label{sec:components}

The B-series papers collectively define fifteen algorithm components, organised
along three orthogonal axes.
Each axis is a dial that can be set independently, and the legal requirements in
Section~\ref{sec:legal-requirements} map cleanly to at most one dial per requirement.
This orthogonality is the key property that makes the Callais disentanglement
constraint satisfiable: the race dial (VRASection) and the partisan dial
(ProportionalSection) are on different axes and can never be set simultaneously
without triggering the \texttt{callais\_preflight} gate.

Figure~\ref{fig:three-axes} provides a visual overview of the three axes and
their components.

\begin{figure}[h]
\centering
\fbox{%
\begin{minipage}{0.9\textwidth}
\centering
\textbf{Three-Axis Algorithm Taxonomy}\\[1em]
\begin{tabular}{p{3.5cm}|p{3.5cm}|p{4.5cm}}
\textbf{Axis A: Vertex weights} & \textbf{Axis B: Edge weights} & \textbf{Axis C: Division strategy} \\[0.5em]
What METIS balances & Boundary cost function & How the bisection tree is structured \\
\hline
\textit{ncon=1, pop only} & \textit{Unweighted} & \textit{Balanced bisection} \\
Standard (B.1, B.7, B.8) & Minimal intervention & Standard $k/2:k/2$ \\[0.5em]
\textit{ncon=2, pop+area} & \textit{Geographic (TIGER)} & \textit{GeoSection (B.8)} \\
AreaSection (B.9) & Compactness & Isoperimetric ratio scan \\[0.5em]
\textit{ncon=1+alignment} & \textit{County-sticky ($\alpha$)} & \textit{AreaSection (B.9)} \\
VRASection (B.14) & Subdivision preservation & Dual-constraint ratio scan \\[0.5em]
\textit{ncon=2, pop+D\_votes} & & \textit{ApportionRegions (B.11)} \\
ProportionalSection (B.12) & & Prime factorisation tree \\[0.5em]
 & & \textit{NestSection (B.13)} \\
 & & Compatible spine \\[0.5em]
 & & \textit{StabilitySection (B.15)} \\
 & & Cross-census CSS filter \\
\end{tabular}
\end{minipage}%
}
\caption{Three-axis taxonomy of the algorithmic redistricting toolbox.
  Axis A controls what METIS optimises (population fractions to balance).
  Axis B controls how edges are weighted (which boundaries are expensive to cut).
  Axis C controls how the bisection tree is structured.
  Each legal requirement maps to at most one component on one axis.}
\label{fig:three-axes}
\end{figure}

\subsection{Axis A: Vertex Weights (What METIS Balances)}

METIS's graph bisection algorithm balances one or two ``constraints'' per vertex,
specified as integer weight vectors.
The choice of what to balance is the first algorithmic design decision.

\subsubsection*{A1: Population Only (\textit{ncon=1})}

The baseline choice: each vertex carries one weight (population from the census),
and METIS balances population fractions.
This is the standard used in B.1 (recursive bisection) and B.7 (solution space
certification).
It satisfies R1 (population equality) and is neutral with respect to all other
requirements.

\textbf{METIS configuration}: $\mathtt{ncon}=1$, $\mathtt{ubvec}[0]=1.001$.

\subsubsection*{A2: Population + Land Area (\textit{ncon=2}) --- AreaSection}

AreaSection (B.9) adds a second METIS constraint: land area.
Each vertex carries two weights (population, area in hectares).
METIS simultaneously targets 50\,\% population balance and 50\,\% area balance
at the first bisection level, with the population constraint taking priority
when the two conflict.

The area constraint eliminates \emph{urban peeling}---the systematic first-level
separation of compact metropolitan cores from the rural remainder.
Urban peeling is not a constitutional violation but creates implicitly packed
maps even with zero partisan input: every recursion level concentrates Democratic
voters into one compact district.
Area balance prevents this by requiring that both halves of the state contain
approximately equal land areas, forcing the algorithm to cross urban-rural
geographic boundaries.

\textbf{METIS configuration}: $\mathtt{ncon}=2$,
$\mathtt{ubvec}=[1.001, 1.10]$,
$\mathtt{tpwgts}=[p_L, 0.5, 1-p_L, 0.5]$ for population fraction $p_L$ at ratio $p_L:1-p_L$.

\textbf{Rodden-null result (B.9)}: AreaSection produces identical seat counts to
GeoSection in 76\,\% of states (25/34 with complete data); the 9 divergent states
differ by exactly 1 seat, with no systematic partisan direction.
The mean proportionality gap in competitive states is $-6.2$\,pp under both
algorithms, confirming geographic voter sorting as the dominant effect.

\subsubsection*{A3: Population + Alignment Signal --- VRASection}

VRASection (B.14) does not add a second METIS constraint.
Instead, it modifies the \emph{ratio selection objective} at each bisection level
by adding a geographic alignment score:
\[
  \text{score}(\text{ratio}) = \mathrm{EC}_\text{norm} - w_\text{vra}
  \cdot A \cdot \max(\mathrm{EC}_\text{norm},\, 1)
\]
where $A = |\text{MVAP}_\text{frac}(L) - \text{MVAP}_\text{frac}(R)|$
is the alignment score (how much minority population is concentrated on one side).

This design explicitly avoids the MetisVra failure mode \citep{b3paper}:
minority VAP fractions (0.0--1.0) and census population counts (10,000--200,000)
differ by $60$--$800\times$, causing numerical degeneration in any METIS
multi-constraint formulation that combines them.
VRASection bypasses this by using $A$ only as a \emph{ratio selector}, not as
a METIS constraint.

\textbf{Key property (Callais compliance)}: The alignment score $A$ uses only
the spatial distribution of minority VAP, not a racial target.
At $w_\text{vra} = 0.40$, compactness is weighted 50\,\% more heavily than
alignment.
The algorithm cannot see how many people of any race are in any district;
it sees only whether one side of a proposed bisection has more minority VAP
\emph{concentration} than the other.

\subsubsection*{A4: Population + Partisan Votes (\textit{ncon=2}) --- ProportionalSection}

ProportionalSection (B.12, planned) is the \emph{Callais-mutex partner} of
VRASection.
It adds partisan vote share (D-vote fraction per tract) as a second METIS balance
constraint, targeting proportional representation.

\textbf{Callais constraint}: ProportionalSection is mutually exclusive with
VRASection.
The \texttt{callais\_preflight} gate in the \texttt{redist} CLI enforces
this mutex: any attempt to run both signals simultaneously is rejected with
an explicit \texttt{[BOUNDARY]} error citing \callais{} p.~36.

\subsection{Axis B: Edge Management}

METIS minimises the total cost of cut edges.
Edge weights determine which boundaries are expensive to cut, thereby shaping
the resulting district geometry.

\subsubsection*{B1: Unweighted Edges}

All edges have weight 1 (or boundary length in metres from TIGER, but not
amplified).
This is the baseline of Paper B.1.
METIS minimises the number of cut adjacency relationships, producing districts
that respect geographic adjacency but do not specifically favour compact shapes.

\subsubsection*{B2: Geographic Edge Weights (TIGER Boundary Lengths)}

Paper B.2 (edge-weighted bisection) demonstrated that weighting edges by their
TIGER boundary length produces a 56\,\% improvement in mean Polsby-Popper
compactness over unweighted bisection.
The mechanism is direct: long borders (between adjacent tracts with extensive
shared boundary) are expensive to cut, so METIS prefers to keep them intact.
Districts therefore tend to follow natural geographic contours (rivers,
ridgelines, coastlines) rather than cutting across them.

\textbf{Compactness result (B.2)}: Mean Polsby-Popper improves from PP~$=0.25$
(unweighted) to PP~$=0.39$ at standard geographic weighting, exceeding the
mean Polsby-Popper of 2020 enacted congressional maps (PP~$=0.32$).

\subsubsection*{B3: County-Sticky Weights --- SubdivisionBisect}

Paper B.10 introduces the governmental stickiness weight:
\[
  w'(u,v) = w_\text{geo}(u,v) \times (1 + \alpha_\text{county} \cdot \mathbf{1}[\text{county}(u) = \text{county}(v)])
\]
Pairs of tracts sharing the same county FIPS code are harder to separate by
a factor of $(1 + \alpha_\text{county})$.
As $\alpha_\text{county}$ increases, county splits are reduced, approached
only when population-equality arithmetic forces them.
The formula generalises to municipality, township, and precinct levels via
$\alpha_\text{mcd}$, $\alpha_\text{place}$, $\alpha_\text{vtd}$.

\textbf{Legal hook}: ``Preserving political subdivisions'' is an explicit
statutory criterion in most state redistricting codes.
B.10 provides the first formal algorithmic operationalisation of this criterion
as a tunable parameter.

\subsection{Axis C: Division Strategy}

The division strategy determines how the recursive bisection tree is structured:
what split ratios are tried, how ratios are selected, and how the tree is
linked to other legal constraints.

\subsubsection*{C1: Balanced Bisection (Standard)}

The standard approach: always split at $\lfloor k/2 \rfloor : \lceil k/2 \rceil$.
This is the baseline of B.1 (recursive bisection) and B.7 (solution space).
It produces a perfectly balanced binary tree, which is computationally efficient
and well-studied.
Its weakness is the caterpillar pathology: in states with compact urban cores,
the balanced split at each level \emph{still} costs less than any alternative
because the urban boundary is shorter than any geographic bisection of the full
state---even after the 1:(k-1) ratio has already been excluded from consideration.

\subsubsection*{C2: GeoSection --- Isoperimetric Ratio Scan}

GeoSection (B.8) scans all feasible split ratios $1:(k-1)$ through
$\lfloor k/2 \rfloor:\lceil k/2 \rceil$ and selects the ratio minimising
$\mathrm{EC}_\text{norm} = \mathrm{EC}/\sqrt{\min(i, k-i)}$.

The $\sqrt{\min(i,k-i)}$ denominator is the isoperimetric correction: for a
convex region enclosing population fraction $i/k$, the minimum boundary length
scales as $\sqrt{\min(i,k-i)/k}$.
This makes all ratios compete on equal footing: the $1:(k-1)$ ratio must win
the isoperimetric competition against the balanced split, not merely be shorter
in absolute terms.

\textbf{Key finding (B.8)}: GeoSection selects the balanced split (or near-equal
split within 20\,\%) for 12 of 44 multi-district states.
For 32 states, it selects asymmetric ratios---but only where the urban core
passes the isoperimetric test.
North Carolina ($k=14$) is the most important case: without normalisation, the
Charlotte/Raleigh peel ($1:13$) wins; with normalisation, the east/west
bisection ($6:8$) wins, reflecting the Piedmont corridor's actual geographic
structure.

\subsubsection*{C3: AreaSection --- Dual-Constraint Ratio Scan}

AreaSection (B.9) combines the GeoSection ratio scan with the dual-constraint
vertex weights of A2.
The Lorenz-curve pre-filter eliminates ratios where simultaneous population-and-area
balance is geometrically infeasible before running any METIS calls.

\textbf{Key finding (B.9)}: AreaSection shifts the natural ratio from asymmetric
to near-equal in many densely urban states: Illinois moves from 1:17 to 8:9,
New York from 1:16 to 13:13, Florida from 1:27 to 14:14.
States with genuinely concentrated populations (Georgia 3:11, Massachusetts 2:7)
retain asymmetric splits because the area constraint cannot be satisfied at
more balanced ratios.

\subsubsection*{C4: ApportionRegions --- Prime Factorisation Tree}

ApportionRegions (B.11, planned) derives the bisection tree structure from the
prime factorisation of $k$, linking it to the Huntington-Hill apportionment
sequence.
For $k = 8 = 2^3$, the tree is a perfect three-level binary tree.
For $k = 14 = 2 \times 7$, the first split is $7:7$, then each half splits
recursively using the factorisation of 7 (prime, so near-equal splits all the
way).

\textbf{Auditability property}: The bisection tree is entirely determined by
$k$, the seed count, and the state's tract adjacency graph.
There are no ``hidden'' parameters.
The plan manifest need only record $k$ and the seed; the full tree follows
deterministically from the published algorithm.

\subsubsection*{C5: NestSection --- Compatible Factorisation Spine}

NestSection (B.13) extends ApportionRegions to multi-chamber consistency.
Given congressional ($C$), senate ($S$), and house ($H$) seat counts,
NestSection computes $g = \gcd(C, S, H)$ and derives a \emph{compatible
factorisation spine} that all three chambers share.

\textbf{Anti-gerrymandering property (B.13)}: A shared geographic hierarchy
is an anti-gerrymandering mechanism by construction.
A partisan actor who locks in a congressional spine cannot draw senate boundaries
that contradict it, because the senate tree is geometrically derived from the
same trunk sequence.
Eleven states achieve score~0 (exact nesting): Oregon, Alabama, Montana, Wyoming,
and others where $g = \min(C, S, H)$.

\subsubsection*{C6: StabilitySection --- Cross-Census CSS Filter}

StabilitySection (B.15) is not a division strategy per se, but a post-hoc
filter that evaluates whether the GeoSection natural ratio is stable across
census years.
The Census Stability Score (CSS) is computed as a weighted combination of
seat stability, ratio stability, and proportionality-gap similarity.

\textbf{Key finding (B.15)}: 67\,\% of comparable states are stable at
$\Delta f \leq 0.05$ across the 2000--2010 census decade.
Iowa is the canonical unstable state: its GeoSection natural ratio shifted from
$1:3$ (asymmetric Des Moines peel) in 2000 to $2:2$ (near-equal east/west split)
in 2010, reflecting dramatic suburban population dispersal.
States with fixed municipal boundaries (Milwaukee, Philadelphia, Chicago) are
highly stable; states with rapidly expanding metro footprints (Phoenix, Las Vegas)
are low-stability.
\end{document}
